\section*{Capítulo \ref{ch:g12:contdist} --- Variables aleatorias continuas}

\begin{solution}{exo:g12:contdist:1}
$\P(X \leq 5) = \frac{5}{15} = \frac13$;
$\P(4 \leq X \leq 10) = \frac{6}{15} = \frac25$;
$\E(X) = \frac{15}{2} = 7.5$ min;
$\sigma(X) = \frac{15}{\sqrt{12}} = \frac{5\sqrt3}{2} \approx 4.3$ min.
\end{solution}

\begin{solution}{exo:g12:contdist:2}
$f \geq 0$ y
$\int_0^2 \frac38 t^2\,\dd t = \frac38\left[\frac{t^3}{3}\right]_0^2
= \frac38 \times \frac83 = 1$: es una \omterm{def:g12:contdist:density}{densidad}.
\[
\E(X) = \int_0^2 \frac38 t^3\,\dd t = \frac38 \times \frac{16}{4} = \frac32,
\qquad
\P(X \geq 1) = \int_1^2 \frac38 t^2\,\dd t = \frac{8 - 1}{8} = \frac78 .
\]
\end{solution}

\begin{solution}{exo:g12:contdist:3}
\emph{1.} $\E(X) = \frac{1}{0.2} = 5$ años;
$\P(X > 5) = \eu^{-0.2\times5} = \eu^{-1} \approx 0.37$.

\emph{2.} Por la \omterm{thm:g12:contdist:memoryless}{falta de memoria} (\cref{thm:g12:contdist:memoryless}),
$\pcond{X > 3}{X > 8} = \P(X > 5) = \eu^{-1} \approx 0.37$: los tres años
de servicio no cambian nada.
\end{solution}

\begin{solution}{exo:g12:contdist:4}
$\P(X > m) = \eu^{-\lambda m} = \frac12$ da
$m = \frac{\ln 2}{\lambda} \approx \frac{0.69}{\lambda}$, menor que
$\E(X) = \frac1\lambda$. La \omterm{def:g12:contdist:density}{densidad} exponencial tiene una cola derecha
larga: unas pocas duraciones de vida atípicamente largas tiran de la
\emph{\omterm{def:g11:stat:mean}{media}} por encima de la \emph{\omterm{def:g11:stat:median}{mediana}}, el valor que supera la mitad
de la población. (Es el periodo de semidesintegración del
\cref{ex:g12:exp:halflife}: la mitad de los átomos lo sobrevive, aunque la
duración de vida \omterm{def:g11:stat:mean}{media} sea mayor.)
\end{solution}

\begin{solution}{exo:g12:contdist:5}
$168 = \mu - \sigma$ y $182 = \mu + \sigma$: en torno al $68\,\%$.
$189 = \mu + 2\sigma$: por encima queda la mitad del
$100 - 95.4 = 4.6\,\%$ restante, es decir, en torno al $2.3\,\%$.
$154 = \mu - 3\sigma$: en torno al $\frac{100 - 99.7}{2} = 0.15\,\%$.
\end{solution}

\begin{solution}{exo:g12:contdist:6}
$\P(X < 500) \leq 2.5\,\%$ significa que $500$ tiene que quedar al menos
dos desviaciones típicas por debajo de la \omterm{def:g11:stat:mean}{media} (la normal deja un
$2.3\,\% \approx 2.5\,\%$ por debajo de $\mu - 2\sigma$; con el $1.96$
convencional, el razonamiento es idéntico):
\[
\mu - 2\sigma \geq 500 \iff \mu \geq 500 + 8 = 508 \text{ g}.
\]
La máquina tiene que ajustarse a $508$ g de \omterm{def:g11:stat:mean}{media}: el precio de la
garantía son $8$ g de producto ``regalado'' por bolsa.
\end{solution}

\begin{solution}{exo:g12:contdist:7}
\emph{1.} Con $p = \frac16$ y $n = 1200$:
$\sqrt{\frac{p(1-p)}{n}} = \sqrt{\frac{5/36}{1200}} \approx 0.0108$, así
que el \omterm{def:g10:numbers:interval}{intervalo} de fluctuación al $95\,\%$ es
\[
\frac16 \pm 1.96 \times 0.0108 \approx \intcc{0.146}{0.188}.
\]

\emph{2.} La frecuencia observada $\frac{260}{1200} \approx 0.217$ queda
muy fuera: la hipótesis de equilibrio se rechaza al nivel del $5\,\%$. La
desviación es de unas $4.6$ desviaciones típicas: para una \omterm{def:g10:numbers:approx}{aproximación}
normal, una \omterm{def:g10:proba:distribution}{probabilidad} del orden de $10^{-6}$, mucho más concluyente que
la cota $\leq 4.6\,\%$ de Bienaymé--Chebyshev (\cref{exo:g12:sums:7}).
\end{solution}

\begin{solution}{exo:g12:contdist:8}
\emph{1.} El margen $\frac{1}{\sqrt n} \leq 0.02$ exige $n \geq 2500$
personas.

\emph{2.} Para distinguir resultados separados por $1$ punto, el margen
tiene que quedar bastante por debajo de $0.5$ puntos, lo que exige
$n \geq \left(\frac{1}{0.005}\right)^2 = 40\,000$ según la fórmula
simplificada: ya inviable para la mayoría de las encuestas. Peor aún: el
margen estadístico solo tiene en cuenta el error de \emph{muestreo}; los
sesgos sistemáticos (\omterm{def:g10:proba:sample}{muestras} no representativas, personas que no
responden, cambios de última hora) no se encogen al crecer $n$ y suelen
superar el punto. Una carrera separada por $1$ punto está, de verdad,
demasiado igualada para pronunciarse.
\end{solution}

\begin{solution}{exo:g12:contdist:9}
\emph{1.} Como $\alpha > 0$, $c \leq \alpha X + \beta \leq d
\iff \frac{c - \beta}{\alpha} \leq X \leq \frac{d - \beta}{\alpha}$, luego
\[
\P(c \leq Y \leq d)
= \int_{(c-\beta)/\alpha}^{(d-\beta)/\alpha} f(t)\,\dd t .
\]
El cambio de variable $y = \alpha t + \beta$ (es decir, leer el área en la
variable $y$, con $t = \frac{y - \beta}{\alpha}$ y
$\dd t = \frac{\dd y}{\alpha}$) lo convierte en
$\int_c^d \frac{1}{\alpha} f\left(\frac{y-\beta}{\alpha}\right)\dd y$: la
\omterm{def:g11:func:function}{función} $g(y) = \frac1\alpha f\left(\frac{y-\beta}{\alpha}\right)$ es una
\omterm{def:g12:contdist:density}{densidad} de $Y$.

\emph{2.} Con $f = \varphi$, $\alpha = \sigma$ y $\beta = \mu$:
\[
g(y) = \frac{1}{\sigma}\,\varphi\!\left(\frac{y - \mu}{\sigma}\right)
= \frac{1}{\sigma\sqrt{2\pi}}\,
\exp\left(-\frac{(y-\mu)^2}{2\sigma^2}\right),
\]
que es, por tanto, la \omterm{def:g12:contdist:density}{densidad} de
$\mathcal N(\mu, \sigma^2) = \mu + \sigma\,\mathcal N(0,1)$.
\end{solution}

\begin{solution}{pb:g12:contdist:1}
\textbf{1.} $\P(\text{espera} \leq 10) = \frac{10}{30} = \frac13$;
$\E = 15$ min; $\sigma = \frac{30}{\sqrt{12}} \approx 8.7$ min.

\textbf{2.} $\int_0^1 2x\,\dd x = 1$: es una \omterm{def:g12:contdist:density}{densidad}.
$\E(X) = \int_0^1 2x^2\,\dd x = \frac23$;
$\P\left(X > \frac12\right) = \int_{1/2}^1 2x\,\dd x = 1 - \frac14 =
\frac34$.

\textbf{3.} $\P(T > 15) = \eu^{-15/10} \approx 0.22$. \omterm{def:g11:stat:median}{Mediana}:
$\eu^{-m/10} = \frac12$: $m = 10\ln 2 \approx 6.9$ min, menor que la \omterm{def:g11:stat:mean}{media}
$10$ porque la larga cola derecha de la exponencial (esperas enormes y
raras) tira de la \omterm{def:g11:stat:mean}{media} hacia arriba mientras deja la \omterm{def:g11:stat:median}{mediana} en la espera
``típica''.

\textbf{4.} Por la \omterm{thm:g12:contdist:memoryless}{falta de memoria},
$\P(T > 35 \mid T > 20) = \P(T > 15) \approx 0.22$: los veinte minutos ya
cumplidos no compran nada; el flujo no envejece.

\textbf{5.} (a) exponencial (la desintegración radiactiva es el modelo para
el que se inventó la \omterm{thm:g12:contdist:memoryless}{falta de memoria}); (b) uniforme sobre $\intcc{0}{8}$
(un horario con un desfase aleatorio); (c) exponencial; (d) ninguno de los
dos: una bombilla desgastada tiene \emph{más} \omterm{def:g10:proba:distribution}{probabilidad} de fallar en la
próxima hora que una nueva, así que
$\P(T > s + t \mid T > s) < \P(T > t)$: el envejecimiento contradice el
teorema de la \omterm{thm:g12:contdist:memoryless}{falta de memoria}.

\textbf{6.} $68\,\%$, $95\,\%$ y $99.7\,\%$: los tres \omterm{def:g10:numbers:interval}{intervalos} de
referencia.

\textbf{7.} $z = \frac{190 - 172}{8} = 2.25$: cola superior
$\approx 1.2\,\%$.

\textbf{8.} $z = 2$: en torno al $2.3\,\%$, aproximadamente una persona de
cada $44$.

\textbf{9.} La especificación está a $\pm 2\sigma$: pasa en torno al
$95.4\,\%$ y se rechaza el $4.6\,\%$, ruinoso a gran escala. A
$\pm 6\sigma$, la tasa de fallo cae a unas dos partes por \emph{mil
millones}: el seis sigma compra el derecho a producir en masa sin
inspeccionar en masa.

\textbf{10.} $\sigma = 5$; con corrección por \omterm{def:g12:limcont:continuity}{continuidad},
$\P \approx \P\left(\abs Z \leq \frac{5.5}{5}\right) =
\P(\abs Z \leq 1.1) \approx 0.73$. El teorema alisa la tabla de Galton del
\cref{pb:g11:binom:1}: la escalera de casillas se convierte en la campana
\omterm{def:g12:limcont:continuity}{continua}.

\textbf{11.} $n \approx \frac{1.96^2 \times 0.25}{0.005^2} \approx
38\,400$ personas; y con ese tamaño, los errores \emph{ajenos} al muestreo
(el marco, las negativas a responder, las mentiras) aplastan al margen:
una carrera de un punto queda fuera del alcance de una encuesta honrada, y
por eso los institutos serios dicen entonces ``demasiado igualado para
pronunciarse''.

\textbf{12.} \omterm{thm:g12:contdist:memoryless}{Falta de memoria}: desde el instante en que llegas, la espera
restante es exponencial de \omterm{def:g11:stat:mean}{media} $10$; los $10$ minutos enteros, no $5$. El
``\omterm{def:g10:numbers:interval}{intervalo} medio de $10$'' del horario induce a error: tu espera tiene la
misma \omterm{def:g12:randvar:rv}{distribución} que un \omterm{def:g10:numbers:interval}{intervalo} completo.

\textbf{13.} Las llegadas al azar caen proporcionalmente a la longitud del
\omterm{def:g10:numbers:interval}{intervalo}: $\P(\text{intervalo largo}) = \frac{15}{20} = \frac34$. Espera
esperada: $\frac34 \times \frac{15}{2} + \frac14 \times \frac52 =
5.625 + 0.625 = 6.25$ min, más que el ingenuo $5$, aunque el \omterm{def:g10:numbers:interval}{intervalo}
medio sea $10$. El culpable: el \emph{muestreo sesgado por la longitud};
los \omterm{def:g10:numbers:interval}{intervalos} largos atrapan a más pasajeros.

\textbf{14.} Las esperas hacia atrás y hacia delante son ambas
exponenciales de \omterm{def:g11:stat:mean}{media} $10$ (la \omterm{thm:g12:contdist:memoryless}{falta de memoria} funciona en los dos
sentidos desde tu llegada): el \omterm{def:g10:numbers:interval}{intervalo} que te contiene mide, en promedio,
$20$ minutos, el doble del \omterm{def:g10:numbers:interval}{intervalo} típico. La paradoja de la inspección
en una frase: \emph{el \omterm{def:g10:numbers:interval}{intervalo} que te toca inspeccionar no es un
\omterm{def:g10:numbers:interval}{intervalo} típico, porque tenías más \omterm{def:g10:proba:distribution}{probabilidad} de caer en uno grande.}

\textbf{15.} Clase \omterm{def:g11:stat:mean}{media}: $\frac{9 \times 10 + 110}{10} = 20$ estudiantes.
\omterm{def:g11:stat:mean}{Media} vivida por el estudiante:
$\frac{90 \times 10 + 110 \times 110}{200} = \frac{13\,000}{200} = 65$
estudiantes. El folleto imprime $20$; tres de cada cinco estudiantes se
sientan en la clase de $110$ y viven el $65$.

\textbf{16.} Las personas populares aparecen en muchas listas de amigos, así
que muestrear ``un amigo'' está sesgado hacia los sociables: tus amigos
salen de los \omterm{def:g10:numbers:interval}{intervalos} largos del horario social.

\textbf{17.} Para $t \geq 0$:
$\P(T > t) = \P(-10\ln U > t) = \P\left(U < \eu^{-t/10}\right) =
\eu^{-t/10}$: exactamente la \omterm{def:g11:func:function}{función} de supervivencia de la exponencial;
un solo \omterm{def:g12:exp:ln}{logaritmo} convierte el ruido uniforme del ordenador en cualquier
tiempo de espera.

\textbf{18.} Cada uniforme tiene \omterm{def:g11:stat:mean}{media} $\frac12$ y \omterm{def:g12:randvar:exp}{varianza}
$\frac{1}{12}$: la suma de doce tiene \omterm{def:g11:stat:mean}{media} $6$ y \omterm{def:g12:randvar:exp}{varianza} $1$, así que la
receta devuelve \omterm{def:g11:stat:mean}{media} $0$ y \omterm{def:g12:randvar:exp}{varianza} $1$. Doce empujoncitos \omterm{def:g12:sums:indep}{independientes}
sumados: el mecanismo de Galton, y con la bendición de De Moivre--Laplace
el histograma ya tiene forma de campana a simple vista.

\textbf{19.} Sobre el modelo: los mercados no son sumas de muchas causas
pequeñas \emph{\omterm{def:g12:sums:indep}{independientes}}; los pánicos lo correlacionan todo (la
lección de 2008 del problema anterior) y producen colas gruesas que ninguna
\omterm{def:g12:contdist:density}{densidad} normal posee. Cuando los datos susurran $10^{-137}$, el catecismo
del estadístico responde: \emph{lo que se rechaza es el modelo, no el
día}.

\textbf{20.} Uniforme: todas las posiciones por igual, el honrado ``solo
conozco los límites''. Exponencial: el riesgo que nunca envejece;
desintegraciones, llegadas, chasquidos. Normal: la campana democrática de
muchas causas pequeñas e \omterm{def:g12:sums:indep}{independientes}; estaturas, errores, promedios. El
aguijón: lo que muestreas está sesgado por la manera en que te tropezaste
con ello; autobuses, clases, amigos. Y el arco: la criatura que contaba
cartones de zumo en el año 6 ha medido, diez volúmenes de problemas
después, a la multitud entera con una curva; el épsilon, la \omterm{def:g12:integ:area}{integral} y el
teorema central del límite esperan en los volúmenes universitarios para
explicar \emph{por qué} la campana dobla por todo.
\end{solution}
